% =====================================================================
% Text to add to the MLmrpi paper to announce the Python package
% Repository: https://github.com/CivilEngineeringSoft/MLmrpi (version 1.1.0)
% Add the Zenodo DOI once the release is archived.
% =====================================================================

% ---------------------------------------------------------------------
% (1) ABSTRACT — add at the end of the abstract
% ---------------------------------------------------------------------
To facilitate its adoption, the proposed framework is released as an open-source
Python package, \texttt{mlmrpi}, available at
\url{https://github.com/CivilEngineeringSoft/MLmrpi}.

% ---------------------------------------------------------------------
% (2) INTRODUCTION — add to the list of contributions (last paragraph)
% ---------------------------------------------------------------------
Finally, to make the $MLmrpi$ readily accessible to researchers and practitioners,
the complete ranking and validation procedure is implemented in an open-source
Python package, \texttt{mlmrpi}, which computes the index from any table of
precomputed performance metrics in a single function call.

% ---------------------------------------------------------------------
% (3) METHODOLOGY — new subsection, to be placed after Method V (MRNV),
%     replacing the paragraph "To implement these steps, we used the Python
%     programming language ..." (or right after it)
% ---------------------------------------------------------------------
\subsection{Implementation: the \texttt{mlmrpi} Python package}\label{Methodology:MLmrpi_package}

To make the proposed index accessible beyond this study, the five-step procedure
of Figure~\ref{fig:mlmrpi_methodology} and the validation analyses of
Sections~\ref{subsection:MLmrpi_robustness} and~\ref{Methodology:Ranking_Agreement}
were implemented in an open-source Python package, \texttt{mlmrpi}, in the spirit
of similar tools developed for classification problems \citep{Akshay2023}. The
package relies only on standard scientific libraries (numpy, pandas, scipy and
matplotlib). It was developed at the Université Claude Bernard Lyon 1 (LMC$^2$,
EA~7427) and is distributed under the MIT licence through a public GitHub
repository (\url{https://github.com/CivilEngineeringSoft/MLmrpi}), from which it can be
installed with a single command:

\begin{verbatim}
pip install git+https://github.com/CivilEngineeringSoft/MLmrpi.git
\end{verbatim}

The input of the package is the performance-metric table described in Step~1
(one row per ML model, one column per precomputed metric, e.g. $R^2$, MSE, RMSE,
MAE or training time for the training, testing and/or complete datasets), supplied
as a pandas DataFrame or as a CSV/Excel file. The optimization direction of each
metric (larger-the-better or smaller-the-better) is inferred automatically from its
name and can be overridden by the user. The package then performs Steps~2 to 5 for
the AW, DEAR, GRA, FA and MRNV methods and returns, for each model, the $MLmrpi$
value, the rank obtained with each method and a consensus rank. The GRA
distinguishing coefficient ($\lambda=0.5$ by default), the eigenvalue threshold of
the factor analysis (Kaiser criterion, eigenvalue $>1$) and optional metric weights
are exposed as parameters.

Besides the ranking itself, the package provides the tools used in this work to
verify the reliability of a ranking: the classical rank-sum aggregation of the
individual metrics; Cohen's $\kappa$, Fleiss' $\kappa$ and Kendall's $W$ for
the agreement among methods and metrics; Pearson and Spearman correlations between
$MLmrpi$ and the individual metrics; the Top (25\,\%), Middle (50\,\%) and Bottom
(25\,\%) performance grouping;
and the Shapiro--Wilk and Wilcoxon signed-rank tests. Results can be exported to
an Excel workbook and visualized with the same figure types and colours as in
this paper (grouped $MLmrpi$ bar charts with ranks, $\kappa$ and correlation heat
maps, Fleiss' $\kappa$ and Kendall's $W$ bar charts). A minimal example is given below:

\begin{verbatim}
import pandas as pd
from mlmrpi import MLmrpi

table = pd.read_csv("metrics.csv", index_col="Model")
r = MLmrpi().fit(table)       # AW, DEAR, GRA, FA and MRNV
print(r.summary())            # MLmrpi values and ranks
print(r.agreement())          # Cohen's kappa, Fleiss' kappa, Kendall's W
r.plot("bars")                # ranking chart (as in Figure 11)
\end{verbatim}

The same analysis is available from the command line
(\texttt{mlmrpi rank metrics.csv}). When the table contains the metrics of several
data partitions (training, testing and complete dataset), the \texttt{PartitionAnalysis}
class (or \texttt{mlmrpi analyze metrics.csv}) ranks each partition and the joint
training--testing set, and generates all the tables and figures presented in
Section~\ref{subsec:MLmrpi_Rankings} with the same layout and colours. The repository also contains the three
metric tables analysed in this paper (Tables~\ref{tab:Performance_ML_models},
\ref{Case_study1_data} and~\ref{Case_study2_data}), a script reproducing the
reported rankings and agreement statistics, and automated unit tests. Because the
package operates only on precomputed metrics, it is independent of the library used
to train the models and applies equally to regression and classification problems.

% ---------------------------------------------------------------------
% (4) CONCLUSION — add as the last point / sentence
% ---------------------------------------------------------------------
The $MLmrpi$ framework has been released as the open-source Python package
\texttt{mlmrpi}, allowing researchers to rank their own models and to assess the
reliability of the resulting rankings with a few lines of code, thereby promoting
transparent and reproducible model selection.

% ---------------------------------------------------------------------
% (5) CODE AVAILABILITY statement (before the references)
% ---------------------------------------------------------------------
\section*{Code availability}
The \texttt{mlmrpi} Python package (version 1.1.0), including the datasets of
Tables~\ref{tab:Performance_ML_models}, \ref{Case_study1_data} and~\ref{Case_study2_data}
and the scripts reproducing the results of this paper, is freely available at
\url{https://github.com/CivilEngineeringSoft/MLmrpi} under the MIT licence (\copyright~2026 Université Claude Bernard Lyon 1).
% Recommended: archive the release on Zenodo and cite its DOI here, e.g.
% (archived version: \url{https://doi.org/10.5281/zenodo.XXXXXXX}).

% ---------------------------------------------------------------------
% (6) BibTeX entry for the reference cited above
% ---------------------------------------------------------------------
% @article{Akshay2023,
%   author  = {Akshay, Akshay and Abedi, Masoud and Shekarchizadeh, Navid and
%              Burkhard, Fiona C. and Katoch, Mitali and Bigger-Allen, Alex and
%              Adam, Rosalyn M. and Monastyrskaya, Katia and Hashemi Gheinani, Ali},
%   title   = {{MLcps}: machine learning cumulative performance score for
%              classification problems},
%   journal = {GigaScience},
%   volume  = {12},
%   pages   = {1--9},
%   year    = {2023},
%   doi     = {10.1093/gigascience/giad108}
% }
